% =====================================================================
% Group 1 synthesis paper: operator algebras / noncommutative geometry
% arc of the GeoVac framework.
%
% Reading guide and unified narrative for Papers 29, 32, 38, 39, 40,
% 42, 43, 44, 45. Math.OA style. Audience: working mathematician in
% NCG or quantum metric geometry considering the GeoVac corpus.
%
% Not a survey of GeoVac as a whole; focused on the operator-algebra
% arc. Assumes background in spectral triples, Lipschitz seminorms,
% Latremoliere propinquity, and Connes-van Suijlekom 2021.
%
% Status: synthesis paper, no new theorems. May 2026.
% =====================================================================
\documentclass[11pt,a4paper]{article}

\usepackage[T1]{fontenc}
\usepackage[utf8]{inputenc}
\usepackage{amsmath, amssymb, amsthm, mathrsfs, mathtools}
\usepackage{geometry}
\geometry{margin=1in}
\usepackage{hyperref}
\hypersetup{colorlinks=true, linkcolor=blue!50!black, citecolor=blue!50!black, urlcolor=blue!50!black}
\usepackage{graphicx}
\usepackage{booktabs}
\usepackage[numbers,sort&compress]{natbib}
\usepackage{xcolor}

% Note: microtype intentionally disabled (matches Papers 38/42/45 to
% avoid the MiKTeX font-expansion environmental issue).

\theoremstyle{plain}
\newtheorem{theorem}{Theorem}[section]
\newtheorem{lemma}[theorem]{Lemma}
\newtheorem{proposition}[theorem]{Proposition}
\newtheorem{corollary}[theorem]{Corollary}
\theoremstyle{definition}
\newtheorem{definition}[theorem]{Definition}
\newtheorem{remark}[theorem]{Remark}
\newtheorem{example}[theorem]{Example}

\newcommand{\nmax}{n_{\max}}
\newcommand{\Nt}{N_{t}}
\newcommand{\sthree}{S^{3}}
\newcommand{\sfive}{S^{5}}
\newcommand{\Hilb}{\mathcal{H}}
\newcommand{\HGV}{\mathcal{H}_{\mathrm{GV}}}
\newcommand{\AGV}{\mathcal{A}_{\mathrm{GV}}}
\newcommand{\DGV}{D_{\mathrm{GV}}}
\newcommand{\DCH}{D_{\mathrm{CH}}}
\newcommand{\DL}{D_{L}}
\newcommand{\JL}{J_{L}}
\newcommand{\Krein}{\mathcal{K}}
\newcommand{\Kplus}{\mathcal{K}^{+}}
\newcommand{\Op}{\mathcal{O}}
\newcommand{\Tcal}{\mathcal{T}}
\newcommand{\Acal}{\mathcal{A}}
\newcommand{\Manifold}{\mathcal{M}}
\newcommand{\opnorm}[1]{\lVert #1 \rVert_{\mathrm{op}}}
\newcommand{\Lipnorm}[1]{\lVert #1 \rVert_{\mathrm{Lip}}}
\newcommand{\cbnorm}[1]{\lVert #1 \rVert_{\mathrm{cb}}}
\newcommand{\norm}[1]{\lVert #1 \rVert}
\newcommand{\R}{\mathbb{R}}
\newcommand{\C}{\mathbb{C}}
\newcommand{\T}{\mathbb{T}}
\newcommand{\N}{\mathbb{N}}
\newcommand{\Z}{\mathbb{Z}}
\newcommand{\Q}{\mathbb{Q}}
\newcommand{\SU}{\mathrm{SU}}
\newcommand{\Uone}{\mathrm{U}(1)}
\newcommand{\SO}{\mathrm{SO}}
\newcommand{\Vol}{\mathrm{Vol}}
\newcommand{\Tr}{\mathrm{Tr}}
\newcommand{\Mat}{M}
\newcommand{\Lprop}{\Lambda^{L}}

\title{The operator-algebra arc of the GeoVac framework:\\
a reading guide for spectral triples, Latr\'emoli\`ere propinquity,
and Lorentzian convergence}
\author{J.~Loutey\thanks{Independent researcher.
\texttt{jloutey@gmail.com}.
This synthesis was prepared in an AI-augmented agentic workflow with
Anthropic's Claude.  All design decisions and editorial choices are
the author's responsibility.}}
\date{May 2026}

\begin{document}
\maketitle

\begin{abstract}
We provide a unified reading of the nine papers that make up the
operator-algebra arc of the GeoVac framework
(Papers~29,~32,~38,~39,~40,~42,~43,~44,~45).  The arc constructs an
explicit almost-commutative spectral triple
$(\AGV, \HGV, \DGV)$ on the Fock-projected unit three-sphere with a
Camporesi--Higuchi spinor Dirac operator, places it in the
Marcolli--van~Suijlekom gauge-network lineage, and proves four
load-bearing theorems.  First, the underlying truncated graphs are
Ramanujan, with integer-algebraic Ihara zeros (Paper~29).  Second, the
Connes--van~Suijlekom truncations $\Tcal_{\nmax}$ converge to the
continuum spectral triple $\Tcal_{\sthree}$ in the Latr\'emoli\`ere
propinquity at quantitative rate
$\Lambda(\Tcal_{\nmax}, \Tcal_{\sthree}) \le C_3 \gamma_{\nmax}$
with $C_3 = 1$ and asymptotic constant $4/\pi$ (Paper~38).
Third, the rate constant $4/\pi$ is universal across the class of
simple compact Lie groups with bi-invariant metric, via a
Plancherel-weight $\times$ Vandermonde-Jacobian cancellation in the
dual-Coxeter normalisation (Paper~40).  Fourth, a $K^{+}$-restricted
weak-form Lorentzian propinquity convergence theorem holds on
truncated $\SU(2)\times\Uone_{t}$ Krein spectral triples (Paper~45) --
to the author's knowledge the first such convergence result in the
math.OA literature, closing a question deferred by Connes and van
Suijlekom \cite{connes_vs2021}.  The constant $4/\pi$ identifies, in
both Riemannian (Paper~38) and Lorentzian (Paper~43~\S 10.2) settings,
as the $M_{1}$ Hopf-base measure signature of the master Mellin
engine catalogued by the case-exhaustion theorem of Paper~32~\S\,VIII.
This synthesis gathers the statements, sketches the load-bearing
proofs, and exhibits the structural relationships between the papers;
no new theorems are introduced.
\end{abstract}

% =====================================================================
\section{Introduction}
\label{sec:intro}
% =====================================================================

The GeoVac framework is a discretization scheme for separable quantum
systems on natural geometries.  When the geometry is the round
$\sthree = \SU(2)$, the framework produces a sequence of finite
truncations whose continuum limit is the Camporesi--Higuchi
spectral triple of the round three-sphere
\cite{camporesi_higuchi1996}.  The present paper synthesizes the nine
documents in the project's operator-algebra group:
Papers~\cite{paper29,paper32,paper38,paper39,paper40_unified,paper42,paper43,paper44,paper45}.
Together they form a self-contained research thread in mathematical
operator algebras with four headline results.

\paragraph{Headline 1: the Hopf graphs are Ramanujan.}
The finite Fock-projected $\sthree$ Coulomb graphs and the
Bargmann--Segal $\sfive$ holomorphic graphs satisfy the
Kotani--Sunada Ramanujan bound at every tested cutoff.  Their Ihara
zeta zeros are algebraic integers with integer coefficients; no $\pi$
or transcendental content appears in the Ihara polynomial
factorizations (Paper~29~\cite{paper29}, Theorems~5.1--5.2 and
Corollary~5.3).

\paragraph{Headline 2: Riemannian propinquity convergence with explicit
rate $4/\pi$.}
The Connes--van~Suijlekom truncations
$\Tcal_{\nmax} = (\Op_{\nmax}, \HGV, \DGV)$ of the
$\SU(2)$ Camporesi--Higuchi spectral triple converge to the continuum
triple $\Tcal_{\sthree}$ in the Latr\'emoli\`ere quantum
Gromov--Hausdorff propinquity~\cite{latremoliere_metric_st_2017,
latremoliere2018}, with
\begin{equation}
\Lambda(\Tcal_{\nmax}, \Tcal_{\sthree})
\;\le\; C_3 \cdot \gamma_{\nmax},
\qquad
C_3 = 1, \qquad
\gamma_{\nmax} \sim \frac{4}{\pi}\,\frac{\log \nmax}{\nmax}
\quad (\nmax \to \infty).
\label{eq:headline_2}
\end{equation}
The asymptotic constant $4/\pi = \Vol(S^{2})/\pi^{2}$ identifies as
the Hopf-base measure of $\sthree \to S^{2}\times S^{1}$ and is
discussed in detail in \S\ref{sec:gh}--\S\ref{sec:mellin} below
(Paper~38~\cite{paper38}, Theorems~3.4 and 4.1).

\paragraph{Headline 3: the rate constant $4/\pi$ is universal across
compact simple Lie groups.}
For any simple compact Lie group $G$ with bi-invariant metric (so
$G \in \{\SU(N), \mathrm{Spin}(n), \mathrm{Sp}(n), G_{2}, F_{4}, E_{6},
E_{7}, E_{8}\}$), the Plancherel-weight $\times$ Vandermonde-Jacobian
cancellation theorem (Paper~40~\S 3.2~\cite{paper40_unified}) leaves
exactly the rank-1 Stein--Weiss constant $4/\pi$ from each positive
root direction's Euler--Catalan asymptote, in the dual-Coxeter
normalisation.  The constant $C_{3} = 1$ remains asymptotically tight
at all ranks via the PRV-summand bound of Kumar (1988) and Vinberg
(1990), applied to the dominant image of the tensor product
$V_{\lambda} \otimes V_{\lambda'}^{*}$.

\paragraph{Headline 4: first Lorentzian propinquity convergence
theorem.}
On truncated $\SU(2)\times \Uone_{T}$ Krein spectral triples
$\Tcal^{L}_{\nmax,\Nt,T}$ built via the Bizi--Brouder--Besnard
\cite{bizi_brouder_besnard2018} signature $(3,1)$ data at
$(m,n) = (4,6)$, the $K^{+}$-restricted weak-form Lorentzian
propinquity satisfies
\begin{equation}
\Lprop(\Tcal^{L}_{\nmax,\Nt,T}, \Tcal^{L}_{\Manifold})
\;\le\; C_3^{\mathrm{joint}} \cdot \gamma^{\mathrm{joint}}_{\nmax,\Nt,T}
\;\longrightarrow\; 0,
\label{eq:headline_4}
\end{equation}
with $\gamma^{\mathrm{joint}}_{\nmax,\Nt,T}
= O(\log \nmax / \nmax + T/\Nt)$, where
$\Lprop(\cdot,\cdot) := \Lambda(\cdot|_{\Kplus},\cdot|_{\Kplus})$ is
the propinquity of the standard metric spectral triple obtained from
the Krein-positive subspace $\Kplus = \{|\psi\rangle : \JL|\psi\rangle
= +|\psi\rangle\}$.  This is, to the author's knowledge, the first
Lorentzian propinquity convergence theorem in the math.OA literature
(Paper~45~\cite{paper45}, Theorem~5.1).  Connes and van~Suijlekom
\cite{connes_vs2021} deferred this question three times to
``elsewhere''; the present arc closes it at the K$^{+}$-weak-form
level on truncated Krein spectral triples.

\paragraph{Position against the literature.}
The arc sits in the lineage of Connes--van~Suijlekom 2021
\cite{connes_vs2021}, which formalised spectral truncations as
operator systems with finite-cutoff Connes distance; of Marcolli and
van~Suijlekom 2014 \cite{marcolli_vs2014}, which introduced
gauge networks as a discrete-graph realisation of finite spectral
triples carrying connection on edges; and of the subsequent
Latr\'emoli\`ere--Hekkelman--McDonald--Toyota--Farsi--Leimbach school
of quantitative propinquity convergence for spectral truncations
\cite{latremoliere_metric_st_2017,
hekkelman2022,hekkelman_mcdonald2024,
hekkelman_mcdonald2024b,toyota2023,
farsi_latremoliere2024,farsi_latremoliere2025,
leimbach_vs2024}.  All these earlier convergence results
are strictly Riemannian (and almost all are on flat structures
$S^{1}$, $\T^{d}$, $S^{2}$ Berezin--Toeplitz).  The Paper~38
$\SU(2)$ result and the Paper~40 generalisation are the first
non-abelian non-flat instances at quantitative rate; the Paper~45
result is the first Lorentzian instance.  Mondino and S\"amann's
synthetic Lorentzian Gromov--Hausdorff program
\cite{mondino_samann2024,mondino_samann2025} is categorically
distinct in object class (pre-length spaces with causal diamonds,
not operator-algebraic spectral triples) and is not in competition
with the present arc.

\paragraph{What this synthesis provides.}
A working mathematician reading the corpus cold may find the
nine-paper sequence difficult to navigate.  This document collects
the constructions, restates the load-bearing theorems with
cross-references back to the original papers, and exhibits the
relationships between them.  No new theorems are introduced.  The
reader is assumed familiar with spectral triples
\cite{connes1995,chamseddine_connes2010}, Lipschitz seminorms and
Rieffel quantum metric spaces, and the Latr\'emoli\`ere propinquity
\cite{latremoliere_metric_st_2017,latremoliere2018}.
Familiarity with \cite{connes_vs2021} is helpful but not strictly
required.  No knowledge of the GeoVac physics-side corpus
(quantum chemistry, QED, precision spectroscopy) is assumed; this
arc is self-contained from the operator-algebra side.

% =====================================================================
\section{The GeoVac spectral triple}
\label{sec:triple}
% =====================================================================

\subsection{The Camporesi--Higuchi data on $\sthree$}

We work throughout on the unit three-sphere $\sthree$, identified
with $\SU(2)$ under its bi-invariant round metric.  Following
Camporesi and Higuchi \cite{camporesi_higuchi1996}, the spinor
bundle is parallelisable, of rank $2$, and the Dirac operator
$\DCH$ has spectrum
\begin{equation}
\mathrm{spec}(\DCH) = \{\pm(n + \tfrac{3}{2}) : n \in \N_{0}\},
\qquad
g^{\mathrm{Dirac}}_{n} = 2(n+1)(n+2)
\label{eq:ch_spectrum}
\end{equation}
on each chirality of the half-integer Hurwitz-shifted shell.  We
write $\Hilb_{\sthree} = L^{2}(\sthree, \Sigma)$ for the chirality-
doubled spinor Hilbert space carrying the Camporesi--Higuchi data.

\subsection{The GeoVac algebra, Hilbert space, and Dirac operator}

The GeoVac construction (Paper~32~\S 3~\cite{paper32}) takes the
following data:

\begin{itemize}\setlength\itemsep{0pt}
\item $\AGV \subset C^{\infty}(\sthree)$ is the involutive algebra
of complex-valued smooth functions on $\sthree$, with the standard
pointwise product and complex conjugation;
\item $\HGV = \Hilb_{\sthree}$, with $\AGV$ acting by pointwise
multiplication;
\item $\DGV = \DCH$, the Camporesi--Higuchi Dirac operator;
\item $J_{\mathrm{GV}}$ is the charge-conjugation real structure
inherited from the spin structure of $\sthree$, satisfying
$J_{\mathrm{GV}}^{2} = -I$ (Euclidean KO-dimension $3$),
$J_{\mathrm{GV}} \DGV = + \DGV J_{\mathrm{GV}}$.
\end{itemize}

The triple $(\AGV, \HGV, \DGV; J_{\mathrm{GV}})$ is real, odd
(no $\Z_{2}$-grading; KO-dimension $3$), and has compact resolvent
$|\DGV|^{-1}$ with eigenvalue multiplicities given by
\eqref{eq:ch_spectrum}.  The Connes axiom audit is recorded in
Paper~32~\S 4~\cite{paper32}, where each of the seven axioms (order
zero, bounded commutators, compact resolvent, order one, real
structure, $\Z_{2}$-grading by convention, and finite summability)
is checked.

\subsection{The Marcolli--van~Suijlekom lineage}

The role of GeoVac in the ecosystem of finite-and-graph spectral
triples is explained in Paper~32~\S 7.  The Marcolli--van~Suijlekom
construction \cite{marcolli_vs2014} attaches a finite spectral
triple to each vertex of a graph and a connection (gauge group
element) to each edge; the spectral action of the resulting object
recovers Wilson lattice gauge theory.  The continuum-limit issue
addressed by Perez-Sanchez \cite{perez_sanchez2024,
perez_sanchez2025} clarifies that the limit is Yang--Mills without
Higgs by default; a Higgs sector appears only via off-diagonal data
on the finite factor.

The GeoVac Fock-projected $\sthree$ graph fits this template
naturally.  Its vertices are the Fock index $(n, \ell, m)$ with
$n = 2j + 1$, $\ell$ the orbital, and $m$ the magnetic quantum
number, providing a Peter--Weyl-indexed basis on $\sthree = \SU(2)$.
Each Fock shell carries a finite spectral triple
(roughly, a finite-dimensional subrepresentation of the
Camporesi--Higuchi triple), and the edge structure encodes the
SO(4) selection rules of the underlying graph.  The graph admits a
$U(1)$ Wilson construction (recorded as GeoVac internal Paper~25),
an $\SU(2)$ Wilson construction on $\sthree = \SU(2)$ (Paper~30),
and an $\SU(3)$ Wilson construction on the Bargmann--Segal $\sfive$
\cite{paper24}, all of which carry the universal kinetic
coefficient $1/(4 N_{c})$.

The Connes axiom audit (Paper~32~\S 4~\cite{paper32}) checks
explicitly that $J^{2} = -I$ and $JD = +DJ$ hold on the truthful
Camporesi--Higuchi data at every finite cutoff
$\nmax \in \{1, 2, 3\}$ tested.  These two relations, together with
the parallelisability of $\sthree = \SU(2)$, distinguish the
GeoVac triple from the Connes--van~Suijlekom circulant comparator,
which is commutative and gives propagation number $1$.

\subsection{Truncations as operator systems}

For each $\nmax \in \N$, the orthogonal projection
$P_{\nmax}: \HGV \to \HGV$ onto the span of Peter--Weyl modes with
$n \le \nmax$ defines a truncated operator system
\begin{equation}
\Op_{\nmax}
\;:=\;
P_{\nmax} \, \AGV \, P_{\nmax}
\;\subset\;
\mathcal{B}(P_{\nmax}\HGV),
\label{eq:operator_system}
\end{equation}
which is self-adjoint, contains the identity, but is in general
\emph{not} multiplicatively closed.  The defect is precisely the
content of the Connes--van~Suijlekom operator-system formalism
\cite{connes_vs2021}: in the natural envelope, $\Op_{\nmax}$ has
propagation number $\mathrm{prop}(\Op_{\nmax}) = 2$ at every tested
$\nmax \in \{2, 3, 4\}$, matching the Toeplitz $S^{1}$
result \cite[Proposition 4.2]{connes_vs2021} verbatim.  A concrete
witness pair at $\nmax = 2$ is $M_{2,1,0}^{2} \notin \Op_{2}$ with
$14.9\%$ Frobenius residual against $\Op_{2}$, deficit concentrated
on the top-shell diagonal in the natural Avery basis
(Paper~32~\S 3.4~\cite{paper32}).  The truncated triple is then
\begin{equation*}
\Tcal_{\nmax} = (\Op_{\nmax}, P_{\nmax}\HGV, P_{\nmax}\DGV P_{\nmax}),
\end{equation*}
with continuum limit
$\Tcal_{\sthree} = (\AGV, \HGV, \DGV)$.

% =====================================================================
\section{The Hopf graphs as Ramanujan substrate}
\label{sec:ramanujan}
% =====================================================================

Before turning to convergence, we record the most basic
spectral-graph property of the underlying finite truncations.

\paragraph{Ihara zeta of the Fock graph.}
For each $\nmax$, let $G_{\nmax}$ denote the directed Fock-projected
$\sthree$ Coulomb graph: vertices are triples $(n, \ell, m)$ with
$n \le \nmax$, $\ell < n$, $|m| \le \ell$; edges encode SO(4)
selection rules.  The Ihara zeta function $\zeta_{G_{\nmax}}(s)$ is
computed in Paper~29~\S 4~\cite{paper29} via the Ihara--Bass
determinantal formula and via the $2E \times 2E$ Hashimoto
non-backtracking edge operator.

\begin{theorem}[Paper~29, Theorems~5.1--5.2]
\label{thm:ramanujan}
For every $\nmax \in \{2, 3, 4, 5\}$, the Fock-projected $\sthree$
graph $G_{\nmax}$ satisfies the Kotani--Sunada Ramanujan bound
strictly:\ the largest non-trivial Hashimoto eigenvalue sits strictly
inside the spectral radius $\sqrt{q_{\max}}$ associated with the
maximum (Perron) vertex degree.  All non-trivial Ihara zeros are
algebraic integers; the corresponding Ihara polynomials factor over
$\Z[s]$ into integer-coefficient pieces, and several factorizations
admit explicit closed forms (e.g.~$(4s^{2}+1)^{2}$, $(9s^{2}+1)^{4}$,
the $12 + 22$ dichotomy for $\sfive$ at $N_{\max} = 3$).
\end{theorem}

\paragraph{Structural consequences.}
Theorem~\ref{thm:ramanujan} secures two facts about the truncated
graphs that are used implicitly elsewhere in the arc.  First, no $\pi$
or transcendental quantity appears anywhere in the Ihara
polynomials, in any of the tested factorizations
(Paper~29~\S 5.4~\cite{paper29}, Corollary on integer-algebraicity).
This is the graph-spectral analog of the
$\pi$-source case-exhaustion theorem (Paper~32~\S 8): both say that
$\pi$ enters only through specific named mechanisms (Hopf-base measure,
spectral action, vertex-parity Hurwitz character), never in the bare
combinatorial data.  Second, the Hashimoto-spectrum Poisson
statistics (Paper~29~\S 6.3) are categorically distinct from the
GUE-like statistics seen on the spectral-zeta side of the framework,
making clear that the Weil dictionary between Selberg and Riemann
breaks at GeoVac by design.

The Ramanujan property of the truncated graphs underwrites our
freedom to take naive sums and traces over $G_{\nmax}$ without
worrying about pathological spectral gaps: the cutoff
graphs are well-conditioned for spectral analysis, and the Ihara
polynomials are entirely in $\Z[s]$.  This is exactly the kind of
substrate one wants under any convergence theorem to follow.

% =====================================================================
\section{Riemannian Gromov--Hausdorff convergence: the central
theorem}
\label{sec:gh}
% =====================================================================

\subsection{Statement (Paper~38)}

The load-bearing theorem of the arc is Paper~38~\cite{paper38},
Theorem~3.4.

\begin{theorem}[Paper~38, GH convergence on $\sthree = \SU(2)$]
\label{thm:gh_su2}
The Connes--van~Suijlekom truncations $\Tcal_{\nmax}$ of the round
$\SU(2)$ Camporesi--Higuchi spectral triple converge to
$\Tcal_{\sthree}$ in the Latr\'emoli\`ere quantum
Gromov--Hausdorff propinquity:
\begin{equation}
\Lambda(\Tcal_{\nmax}, \Tcal_{\sthree})
\;\le\;
C_{3} \cdot \gamma_{\nmax} \;\to\; 0
\quad (\nmax \to \infty),
\label{eq:gh_su2}
\end{equation}
with $C_{3} = 1$ (asymptotic tight via the PRV summand bound) and
asymptotic constant $4/\pi$:
\begin{equation}
n \,\gamma_{n} \;\sim\; (4/\pi)\,\log n
\qquad (n \to \infty).
\label{eq:gh_rate_constant}
\end{equation}
A uniform finite-cutoff bound $\gamma_{n} \le 6\,(\log n)/n$ holds
for all $n \ge 2$.
\end{theorem}

The proof proceeds via five lemmas L1$'$, L2, L3, L4, L5 of
Paper~38~\S 4--\S 6.  We sketch each.

\subsection{Lemma L1$'$: chirality-doubled truncated operator system}

The original Connes--van~Suijlekom truncation is on the scalar Fock
graph; the spinor lift to the full Camporesi--Higuchi bundle (with
both chiralities) is the subject of an internal sprint (R3.5).  The
truthful Camporesi--Higuchi data at finite $\nmax$ has the property
that on most cross-shell pairs the Connes distance is $+\infty$
(an artefact of the $n$-shell degeneracy in the Avery basis); this
is repaired by lifting to the full Dirac sector with both
chiralities and off-diagonal shell-coupling multipliers.  The
chirality-doubled operator system has dimension
$\dim \Op_{\nmax} = 1, 14, 55, 140, \dots$ at
$\nmax = 1, 2, 3, 4, \dots$ and every cross-shell pair has finite
Connes distance.  See Paper~38~\S 4.1~\cite{paper38} for the
explicit construction.

\subsection{Lemma L2: central spectral Fej\'er kernel on $\SU(2)$}

The propinquity tunneling pair built in L5 below requires a positive,
state-preserving, central good kernel on $\SU(2)$ whose
mass-concentration moment vanishes at a known rate.  Paper~38~\S 4.2
constructs the central spectral Fej\'er kernel
\begin{equation}
K_{\nmax}(g) \;=\; \frac{1}{Z_{\nmax}}
\left|\sum_{j \le j_{\max}} \sqrt{2j+1}\,\chi_{j}(g)\right|^{2},
\qquad
j_{\max} = (\nmax-1)/2,
\quad
Z_{\nmax} = \frac{\nmax(\nmax+1)}{2}.
\label{eq:fejer_kernel}
\end{equation}
The Plancherel symbol on the Avery--Wen--Avery basis is
$\widehat{K}(j) = (2j+1)/Z_{\nmax}$ for $j \le j_{\max}$.  The
cb-norm of the associated Herz--Schur multiplier acting on the
group $C^{*}$-algebra is
\begin{equation}
\cbnorm{S_{K_{\nmax}}} \;=\; \frac{2}{\nmax + 1}
\;=\; O\!\bigl(1/\nmax\bigr),
\label{eq:cb_norm}
\end{equation}
via the Bo\.zejko--Fendler central-multiplier equality
\cite{bozejko_fendler1991} on amenable compact groups (Pier
\cite{pier1984}; Paulsen \cite{paulsen2002}; Pisier
\cite{pisier2001}).

\paragraph{The asymptotic constant $4/\pi$.}
The mass-concentration moment $\gamma_{\nmax}$ of the kernel
admits a closed-form decomposition via Stein--Weiss-type integration
by parts \cite{stein_weiss1971}, leading to the
\emph{Stein--Weiss sharpening} (Paper~38, Appendix~A):
\begin{equation}
\gamma_{\nmax}
\;=\;
\pi - \frac{4\, T_{\nmax}}{\pi\, Z_{\nmax}},
\quad
\text{with } \frac{n\,\gamma_{n}}{\log n} \to \frac{4}{\pi}
\;\text{ as } n \to \infty.
\label{eq:gamma_asymptote}
\end{equation}
The constant $4/\pi = \Vol(S^{2})/\pi^{2}$ is the
\emph{Hopf-base measure factor} in the $\SU(2)/\Uone$ Haar
normalisation.  Concretely, for the Hopf bundle
$\sthree \to S^{2}$, $\Vol(S^{2}) = 4\pi$ and the right-hand side of
\eqref{eq:gamma_asymptote} normalises by $\pi^{2}$ from the SU(2)
character integration $\int_{0}^{2\pi} \sin^{2}(\theta/2)\,d\theta$.
The expected one-log slowdown of the non-abelian SU(2) rate compared
with the Leimbach--van~Suijlekom torus rate $O(1/\nmax)$
\cite{leimbach_vs2024} is precisely the $\log n$ factor in
\eqref{eq:gamma_asymptote}.

\subsection{Lemma L3: Lipschitz comparison with $C_{3} = 1$}

For each $f \in C^{\infty}(\sthree)$ and each truncated multiplier
$\Mat_{f}$ on $\Op_{\nmax}$, the Lipschitz comparison
\begin{equation}
\opnorm{[\DCH, \Mat_{f}]}
\;\le\;
C_{3} \cdot \norm{\nabla f}_{L^{\infty}(\sthree)}
\label{eq:lipschitz_bound}
\end{equation}
holds with $C_{3} = 1$ (Paper~38, Theorem~5.2).  The proof uses
SO(4) selection rules: the off-diagonal element of $[\DCH, \Mat_{f}]$
between Peter--Weyl shells $V_{\lambda}$ and $V_{\lambda'}$ is
bounded by $\sqrt{C(\lambda)} - \sqrt{C(\lambda')}$ where
$C(\lambda)$ is the quadratic Casimir of the highest weight; the
Dirac-triangle inequality $|D(\pi) - D(\pi')| \le \sqrt{C(\sigma)}$
on intertwining summands $\sigma \subset \pi \otimes \pi'^{*}$
secures $C_{3} \le 1$, and the Avery harmonic gradient scaling
forces $C_{3} = 1$ in the limit $\nmax \to \infty$.  At any finite
$\nmax$, $C_{3} = (N-1)/\sqrt{N^{2}-1}$ for $N = \nmax$, which
converges to $1$ from below.

The same bound is proved for all simple compact Lie groups in
Paper~40~\S 3.3~\cite{paper40_unified}, with asymptotic tightness
$C_{3}(G) = 1$ at all ranks; see
\S\ref{sec:universality} below.

\subsection{Lemma L4: Berezin reconstruction}

A Berezin map $B_{\nmax}: C(\sthree) \to \Op_{\nmax}$ is constructed
in Paper~38~\S 5.3 as
\begin{equation}
B_{\nmax}(f) \;=\; \sum_{N \le \nmax,\,L,\,M}
\widehat{K}(N)\, c_{NLM}\, M_{NLM},
\quad
\widehat{K}(N) = N / Z_{\nmax}.
\label{eq:berezin}
\end{equation}
This is the $\SU(2)$ non-K\"ahler analog of the Connes--van~Suijlekom
\cite{connes_vs2021} Fej\'er kernel reconstruction on $S^{1}$ and the
non-holomorphic analog of Hawkins \cite{hawkins2000}.  Four properties
are established:
\begin{itemize}\setlength\itemsep{0pt}
\item[(a)] Positivity: for $f \ge 0$, $B_{\nmax}(f)$ is a positive
operator on $P_{\nmax}\HGV$ (via convolution + positivity-preserving
compression);
\item[(b)] Contractivity: $\opnorm{B_{\nmax}(f)} \le \norm{f}_{\infty}$;
\item[(c)] Approximate identity: $\opnorm{B_{\nmax}(f) - P_{\nmax}
M_{f} P_{\nmax}} \le \gamma_{\nmax} \norm{\nabla f}_{\infty}$ with
$\gamma_{\nmax}$ as in L2;
\item[(d)] L3 compatibility: $\opnorm{[\DCH, B_{\nmax}(f)]} \le
\norm{\nabla f}_{\infty}$ (i.e.~$C_{3} = 1$ is inherited via
$\widehat{K} \le 1$).
\end{itemize}
This is the operator-system analog of the Berezin--Toeplitz
quantisation; see Hekkelman--McDonald \cite{hekkelman_mcdonald2024}
for the closely related $\T^{d}$ case.

\subsection{Lemma L5: Latr\'emoli\`ere propinquity assembly}

The tunneling pair $(B_{\nmax}, P_{\nmax})$ between
$\Tcal_{\sthree}$ and $\Tcal_{\nmax}$ has reach
$\le \gamma_{\nmax}$ (by L4(c) approximate identity and L2(g)
Bo\.zejko--Fendler dual estimate), and height contributions are
bounded by L4(b) contractivity and the fact that the projection $P$
contributes $\mathrm{height}_{P} = 0$.  The assembly machinery is
Latr\'emoli\`ere \cite{latremoliere_metric_st_2017}, applied
verbatim.  This gives
\begin{equation}
\Lambda(\Tcal_{\nmax}, \Tcal_{\sthree})
\;\le\; C_{3} \cdot \gamma_{\nmax} \;\to\; 0,
\label{eq:gh_assembly}
\end{equation}
completing the proof of Theorem~\ref{thm:gh_su2}.

\subsection{Numerical verification}

Paper~38 records the explicit $\Lambda$ values at small cutoffs:
\begin{equation}
\Lambda(\Tcal_{2}, \Tcal_{\sthree}) \approx 2.0746,
\quad
\Lambda(\Tcal_{3}, \Tcal_{\sthree}) \approx 1.6101,
\quad
\Lambda(\Tcal_{4}, \Tcal_{\sthree}) \approx 1.3223,
\label{eq:numerical_lambda}
\end{equation}
with $\Lambda(4)/\Lambda(2) = 0.637$.  These are monotone decreasing
and consistent with the $O(\log \nmax / \nmax)$ rate.

\subsection{Tensor extension (Paper~39)}

Paper~39 extends Theorem~\ref{thm:gh_su2} to tensor products
$\Tcal_{\sthree}^{\lambda_{a}} \otimes \Tcal_{\sthree}^{\lambda_{b}}$
of two copies of the same manifold at distinct focal lengths
$\lambda_{a} \ne \lambda_{b}$ \cite{paper39}.  The result is a
proved-at-qualitative-rate theorem closing the published-open NCG
question of convergence for the tensor of two infinite metric
spectral triples.  Two technical points underwrite the closure:
\begin{itemize}\setlength\itemsep{0pt}
\item \textbf{R1: asymmetric Pythagorean Leibniz.}  The Pythagorean
Leibniz bound on tensor commutators (Connes--Marcolli graded form)
combined with an asymmetric-supremum correction gives the full
operator-system constant $C_{3}^{(2)}$ that is bounded away from
$2 \sqrt{2}$ and approaches the single-factor $C_{3} = 1$
asymptotically.
\item \textbf{R2: Latr\'emoli\`ere 2017 explicit
$\varepsilon_{\mathrm{cross}}$.}  The cross term in the tunneling
pair is bounded explicitly by Latr\'emoli\`ere
\cite{latremoliere_metric_st_2017}~\S 4, with rate $O(\max \gamma)$
and constant $\le 2\sqrt{2}$.
\end{itemize}

The bit-identical convergence ratio
$\Lambda(\nmax,\nmax) / \Lambda(2,2)$ on the diagonal matches the
single-factor Paper~38 ratio numerically, providing a load-bearing
sanity check.

% =====================================================================
\section{The universality of the rate constant $4/\pi$}
\label{sec:universality}
% =====================================================================

Paper~40 \cite{paper40_unified} extends Theorem~\ref{thm:gh_su2}
from $\SU(2)$ to the class of all simple compact Lie groups with
bi-invariant metric.

\begin{theorem}[Paper~40, universal rate constant]
\label{thm:universal_rate}
Let $G$ be a simple compact connected Lie group of rank $r \ge 1$
with bi-invariant metric (so
$G \in \{\SU(N), \mathrm{Spin}(n), \mathrm{Sp}(n), G_{2}, F_{4},
E_{6}, E_{7}, E_{8}\}$).  The truncated Camporesi--Higuchi triples
$\Tcal_{\nmax}(G)$ converge to $\Tcal_{G}$ in the Latr\'emoli\`ere
propinquity with the same asymptotic rate constant as
\eqref{eq:gh_rate_constant}: namely $4/\pi$, in the dual-Coxeter
normalisation.  The constant $C_{3}(G) = 1$ remains asymptotically
tight at all ranks.
\end{theorem}

\paragraph{Plancherel-weight $\times$ Vandermonde-Jacobian
cancellation.}
The proof of Theorem~\ref{thm:universal_rate} runs through a
structural identity which we record because it is the deepest
algebraic content of the universality.  In the central spectral
Fej\'er kernel on $G$,
\begin{equation}
K_{\nmax}(g) \;=\; \frac{1}{Z_{\nmax}(G)}
\left|\sum_{\lambda \,:\, |\lambda + \rho|^{2} \le M(\nmax)^{2}}
\sqrt{\dim V_{\lambda}}\,\chi_{\lambda}(g)\right|^{2},
\label{eq:fejer_kernel_general}
\end{equation}
the Plancherel weight $\sqrt{\dim V_{\lambda}}$ on the squared
amplitude (so the weight is $\dim V_{\lambda}$) and the Vandermonde
Jacobian $|\Delta(t)|^{2}$ in the Weyl integration formula on $G$
share the same $\prod_{\alpha > 0}$ product structure over positive
roots.  In the dual-Coxeter normalisation these products cancel
exactly, leaving only the rank-1 Stein--Weiss constant $4/\pi$ from
each positive-root direction's Euler--Catalan asymptote
$\sum_{d \mathrm{ odd}} 1/d^{2} = \pi^{2}/8$.  The universal constant
$4/\pi$ is then \emph{independent} of rank, Weyl group order, and
Lie type.

\paragraph{Asymptotic tightness $C_{3}(G) = 1$ via PRV.}
The Lipschitz comparison bound \eqref{eq:lipschitz_bound} extends to
$G$ via the PRV-summand theorem (Kumar 1988, Vinberg 1990):\ for
every Weyl group element $w \in W$, the dominant image
$\sigma_{w}$ of $\lambda + w\lambda'^{*}$ is a summand of
$V_{\lambda} \otimes V_{\lambda'}^{*}$, and
$|\sigma_{w} + \rho|^{2} \ge |\lambda - \lambda'|^{2} + |\rho|^{2}$.
This covers the asymptotic family
$(\nmax, 0, \dots, 0) \otimes (1, 0, \dots, 0)^{*}$ rigorously, giving
$C_{3}(G) = 1$ asymptotic-tight in all simple Lie algebra types
$A_{n}, B_{n}, C_{n}, D_{n}, G_{2}, F_{4}, E_{6}, E_{7}, E_{8}$.
Interior (non-PRV) summands at finite cutoff are verified
numerically across the panel $\SU(3), \SU(4), \mathrm{Sp}(2), G_{2}$
at $5641/5641$ cells; analytic closure via Brauer--Klimyk signed-sum
cancellation is bookkeeping.

\paragraph{Companion: state-space GH.}
The state-space Gromov--Hausdorff convergence (without rate and
without Dirac) at the same generality was proved by
Gaudillot-Estrada--van~Suijlekom (IMRN 2025, paper rnaf197).
Paper~40 provides the missing rate constant, the Lipschitz spinor
bound, and the asymptotic tightness, lifting Gaudillot-Estrada--vS
from state-space GH to spectral-triple-level propinquity convergence
with explicit rate.

% =====================================================================
\section{Modular structure and Wick rotation}
\label{sec:modular}
% =====================================================================

\subsection{The four-witness Wick-rotation theorem (Paper~42)}

The Wick rotation between Euclidean and Lorentzian field theory is
classically formulated through four equivalent presentations on a
flat spacetime:
\begin{itemize}\setlength\itemsep{0pt}
\item Hartle--Hawking 1976 \cite{hartle_hawking1976}:
the temperature of the Euclidean cigar's $\tau$-circumference
identifies as the Hawking temperature;
\item Sewell 1982 \cite{sewell1982}:
the KMS state at inverse-temperature $\beta = 2\pi/\kappa_{g}$ is the
wedge vacuum;
\item Bisognano--Wichmann 1976 \cite{bisognano_wichmann1976}:
the modular automorphism group of the wedge von~Neumann algebra
acts as boosts around the wedge edge;
\item Unruh 1976 \cite{unruh1976}:
an accelerated observer with proper acceleration $a$ sees a
$T_{U} = a/(2\pi)$ thermal bath in the vacuum.
\end{itemize}
On a flat background, the four are equivalent under standard
Wick-rotation arguments.  On a curved background like $\sthree$ they
are not obviously equivalent, and at finite cutoff (operator-system
level) they require separate construction.

Paper~42 \cite{paper42} establishes the four-witness theorem at the
operator-system level on the truncated Camporesi--Higuchi spectral
triple $\Tcal_{\nmax}$ via two structurally independent but
operator-action conjugate constructions:

\begin{theorem}[Paper~42, BW-$\alpha$ geometric realisation]
\label{thm:bw_alpha}
For each $\nmax \in \{2, 3, 4, 5\}$, there exists a wedge projection
$P_{W} = \tfrac{1}{2}(I + R_{\mathrm{polar}})$ on $P_{\nmax}\HGV$
(bit-exact idempotent and Hermitian) and a geometric modular
Hamiltonian
\begin{equation}
K_{\alpha}^{W} \;=\; J_{\mathrm{polar}}
\quad \text{(rotation generator preserving } W),
\label{eq:k_alpha}
\end{equation}
with integer spectrum $\{\text{two}\, m_{j}\}$ on the full-Dirac basis.
The modular flow satisfies $\sigma_{2\pi}(O) = O$ bit-exact (residual
$O(\sqrt{\dim \Hilb} \cdot \varepsilon_{\mathrm{machine}})$) for the
four witnesses BW, Hartle--Hawking, Sewell, Unruh, with cross-witness
collapse residual identically zero.
\end{theorem}

\begin{theorem}[Paper~42, BW-$\gamma$ Tomita--Takesaki realisation]
\label{thm:bw_gamma}
On the Krein-GNS Hilbert--Schmidt space
$\Hilb_{\mathrm{GNS}} = M_{\dim W}(\C)$ with KMS state
$\rho_{W} = e^{-K_{\alpha}^{W}}/Z$ and the BW choice $H_{\mathrm{local}}
:= K_{\alpha}^{W} / \beta$, the polar decomposition
$S = J_{\mathrm{TT}} \cdot \Delta^{1/2}$ gives a Tomita-Takesaki
modular Hamiltonian $K_{\mathrm{TT}} = -\log \Delta$ with the
same integer spectrum, and the operator-action conjugacy
\begin{equation}
\sigma_{t}^{\mathrm{TT}}(a) \;=\; \sigma_{-t}^{\alpha}(a)
\label{eq:bw_alpha_gamma_conjugacy}
\end{equation}
holds bit-exact at every tested $t \in \{1, \pi, 2\pi\}$.  The
modular operator $\Delta$ satisfies $J_{\mathrm{TT}}^{2} = +I$,
categorically distinct from the Connes $J_{\mathrm{GV}}$ (KO-dim $3$,
$J^{2} = -I$).
\end{theorem}

\paragraph{Load-bearing scope finding (Paper~42~\S 7.2).}
A key derived structural finding is that the local Hamiltonian
required by Tomita--Takesaki at $\beta = 2\pi$ \emph{must} be
$H_{\mathrm{local}} = K_{\alpha}^{W} / \beta$ rather than the
intrinsic Camporesi--Higuchi Dirac operator $\DCH$.  The intrinsic
$\DCH$ has half-integer spectrum, so $\beta \cdot \DCH$ would not have
integer eigenvalues and the bit-exact closure $e^{i \cdot 2\pi \cdot
n} = 1$ would fail.  Spectral-action (Chamseddine--Connes
\cite{chamseddine_connes2010}) and thermal-time (Connes--Rovelli)
paradigms point to different generators on the wedge KMS state.
This is the open question O3 of Paper~42~\S 10, sharpened in
Paper~43 to the signature-independence statement below.

\subsection{Lorentzian extension at finite cutoff (Paper~43)}

Paper~43 \cite{paper43} extends the four-witness theorem to the
Lorentzian Krein space at signature $(3,1)$, built via the
Bizi--Brouder--Besnard prescription \cite{bizi_brouder_besnard2018}
at $(m,n) = (4,6)$ in the Peskin--Schroeder chiral basis, with
Lorentzian Dirac
\begin{equation}
\DL \;=\; i\,\bigl(\gamma^{0} \otimes \partial_{t}
\;+\; \DGV \otimes I_{\Nt}\bigr)
\label{eq:lorentzian_dirac}
\end{equation}
per van~den~Dungen 2016 Proposition 4.1 \cite{vandungen2016}.

\begin{theorem}[Paper~43~\S 5--\S 7, Lorentzian four-witness
at finite cutoff]
\label{thm:lorentzian_four_witness}
At every tested $(\nmax, \Nt) \in \{1, 2, 3\} \times \{1, 11, 21\}$,
the Krein-level analogues of \eqref{eq:k_alpha} and
\eqref{eq:bw_alpha_gamma_conjugacy} hold bit-exact:
\begin{itemize}\setlength\itemsep{0pt}
\item[(i)] The geometric Lorentzian modular Hamiltonian
$K_{L}^{\alpha,W} = J_{\mathrm{polar}}|_{W_{L}}$ has integer
spectrum on the Krein hemispheric wedge.
\item[(ii)] The Krein-Tomita modular operator $\Delta_{L}$ admits a
polar decomposition $S = J_{L,\mathrm{TT}} \cdot \Delta_{L}^{1/2}$;
the Tomita modular Hamiltonian $K_{L}^{\mathrm{TT}}$ has integer
spectrum.
\item[(iii)] $\sigma_{2\pi}^{L,\alpha}(O) = O$ and
$\sigma_{2\pi}^{L,\mathrm{TT}}(a) = a$ bit-exact, with cross-witness
collapse residual zero.
\item[(iv)] At $\Nt = 1$, the Lorentzian construction reduces to the
Riemannian one of Theorems~\ref{thm:bw_alpha}--\ref{thm:bw_gamma}
bit-identically.
\end{itemize}
\end{theorem}

\paragraph{Signature-independence of the load-bearing scope.}
The deepest structural finding of Paper~43 is that the
``$H_{\mathrm{local}} \ne D_{W}^{L}$'' statement is
\emph{signature-independent} at the Riemannian reduction.  At
$\Nt = 1$ the residual norm $\|H_{\mathrm{local}} - D_{W}^{L}\|_{F}$
on the Lorentzian Krein wedge equals the Riemannian-side value at
every $\nmax$ (Paper~43~\S 7.1):
\begin{equation*}
\|H_{\mathrm{local}}|_{(3,1),\Nt=1} - D_{L}^{W}|_{\Nt=1}\|_{F}
\;=\;
\|H_{\mathrm{local}}^{\mathrm{Riem}} - D_{W}\|_{F}
\qquad
(\text{at every } \nmax).
\end{equation*}
At $\Nt > 1$ the Lorentzian residual is refined upward by the
temporal-derivative content of $\DL$: the spectral-action object
$\DL$ has more content than $\DGV$, so the gap between
$H_{\mathrm{local}}$ (modular-Hamiltonian generator) and $D_{W}^{L}$
(spectral-action object) widens with temporal cutoff.  This refines
the Paper~42~\S 10 open question O3 from
``open about signature-dependence'' to a deeper structural feature.

\subsection{Pythagorean HS-orthogonality and the $1/\pi^{2}$
signature}

Paper~43~\S 10.2 records a Pythagorean refinement of the modular vs
spectral-action gap.  Let $\langle \cdot, \cdot \rangle_{\mathrm{HS}}$
denote the Hilbert--Schmidt inner product on the wedge operator
algebra.

\begin{corollary}[Paper~43~\S 10.2, HS-orthogonality on the
Lorentzian wedge]
\label{cor:hs_orthogonality}
For every $\kappa_{g} > 0$ and every panel cell
$\nmax \in \{1, \dots, 6\}$, $\Nt \in \{1, 11\}$, and witness
$w \in \{\mathrm{BW}, \mathrm{HH}_{1}, \mathrm{HH}_{2},
\mathrm{Sew}, \mathrm{Unruh}_{1}, \mathrm{Unruh}_{2}\}$:
\begin{equation}
\langle H_{\mathrm{local}}, D_{W}^{L} \rangle_{\mathrm{HS}}
\;=\; 0
\qquad (\text{bit-exact}),
\label{eq:hs_orthogonal}
\end{equation}
and the residual norm has the closed form
\begin{equation}
\|H_{\mathrm{local}} - D_{W}^{L}\|_{F}^{2}
\;=\;
\frac{\kappa_{g}^{2}}{4\pi^{2}}\,S(n)
\;+\; D(n),
\quad
S(n) = \frac{n(n+1)(n+2)(2n^{2}+4n-1)}{15},
\quad
D(n) = \frac{n(n+1)(n+2)(2n+1)(2n+3)}{20}.
\label{eq:hs_closed_form}
\end{equation}
\end{corollary}

The closed form \eqref{eq:hs_closed_form} is PSLQ-verified at
$100$ digits with coefficient ceiling $10^{6}$ at $18$ panel cells.
The $1/\pi^{2}$ prefactor on $\|H_{\mathrm{local}}\|^{2}$ is the same
$M_{1}$ Hopf-base measure signature
\eqref{eq:gamma_asymptote} that appears in the Paper~38 GH
convergence rate constant $4/\pi$.  This crosswiring between
sections \ref{sec:gh} and \ref{sec:modular} is taken up systematically
in \S\ref{sec:mellin}.

\paragraph{Subspace decomposition.}
The mechanism is the orthogonal decomposition of the operator
algebra $B(\Krein_{W})$ into a diagonal subspace (where
$H_{\mathrm{local}} = K_{\alpha}^{W}/\beta$ lives, since
$K_{\alpha}^{W}$ has eigenvalues $\text{two}\,m_{j}$ on the full-Dirac
basis) and an off-diagonal subspace with $\Delta n = \pm 1$ (where
$D_{W}^{L}$ lives, since the Camporesi--Higuchi Dirac intertwines
$\pm m_{j}$ chirality partners at adjacent shells).  Formal
operator-theoretic proof of this decomposition is named as
Paper~43 open question O4; the PSLQ-verified Pythagorean closed form
\eqref{eq:hs_closed_form} is empirically sharp at $100$ digits
and $18$ panel cells.

% =====================================================================
\section{Lorentzian propinquity}
\label{sec:lorentzian}
% =====================================================================

\subsection{Operator-system substrate (Paper~44)}

The first step in any Lorentzian propinquity proof is an operator-
system substrate.  Paper~44 \cite{paper44} provides it on the
chirality-doubled Camporesi--Higuchi spinor space
$\Krein_{\nmax,\Nt} = \HGV^{\nmax} \otimes \C^{\Nt}$ at BBB signature
$(m,n) = (4,6)$.

\begin{theorem}[Paper~44, propagation number and envelope dependence]
\label{thm:p44_prop_number}
The truncated Lorentzian operator system $\Op^{L}_{\nmax,\Nt}$
satisfies
\begin{equation}
\mathrm{prop}(\Op^{L}_{\nmax,\Nt}) \;=\; 2
\quad \text{under the Weyl-doubled achievable envelope},
\quad
\mathrm{prop}(\Op^{L}_{\nmax,\Nt}) \;=\; \infty
\quad \text{under the full } (\dim \Krein)^{2} \text{ envelope},
\label{eq:p44_prop}
\end{equation}
generalising the Riemannian $\mathrm{prop} = 2$ result for
$\Op_{\nmax}$ (Paper~32~\S 3.4) to the Lorentzian Krein setting.
The witness pair $(M^{2,1,0,0}, M^{2,1,0,0*})$ at $\nmax = 2$
exhibits $\sim 38\%$ Frobenius residual against $\Op^{L}_{2,3}$.
Riemannian-limit recovery at $\Nt = 1$ is bit-exact (Frobenius
residual $= 0.0$ in float64) across $\nmax \in \{1, 2, 3\}$.
\end{theorem}

\paragraph{$K^{+}$-positivity is trivial at the operator-multiplier
level.}
A key structural observation of Paper~44 is that all chirality-
doubled scalar multipliers $M \oplus M$ commute with $\JL =$
chirality-swap, and therefore preserve $\Kplus$ automatically:
$\Op^{L,+} = \Op^{L}$ exactly.  Consequently the non-trivial
Krein-positivity program shifts from operator level to \emph{state}
level (Wasserstein--Kantorovich on $K^{+}$ states), which is the
natural setting for the Paper~45 propinquity proof below.

\subsection{$K^{+}$-restricted weak-form Lorentzian propinquity
convergence (Paper~45)}

Paper~45 \cite{paper45} closes the convergence theorem on the
substrate of Paper~44.

\begin{definition}[Paper~45~\S 1.3, K$^{+}$-weak-form Lorentzian
propinquity]
\label{def:k_plus_prop}
For two truncated Lorentzian spectral triples
$\Tcal_{i}^{L}$ ($i = 1, 2$), define
\begin{equation}
\Lprop(\Tcal_{1}^{L}, \Tcal_{2}^{L})
\;:=\;
\Lambda(\Tcal_{1}^{L}|_{\Kplus}, \Tcal_{2}^{L}|_{\Kplus}),
\label{eq:k_plus_prop_def}
\end{equation}
where $\Lambda$ is the standard Latr\'emoli\`ere propinquity
\cite{latremoliere_metric_st_2017,latremoliere2018} and the
restriction to $\Kplus = \{|\psi\rangle : \JL|\psi\rangle =
+|\psi\rangle\}$ is the standard metric spectral triple obtained from
the Krein-positive subspace, on which the indefinite Krein product
reduces to a positive-definite Hilbert inner product.
\end{definition}

\begin{theorem}[Paper~45, Theorem~5.1]
\label{thm:p45_main}
For the truncated $\SU(2) \times \Uone_{T}$ Krein spectral triples
$\Tcal^{L}_{\nmax,\Nt,T}$, with $\DL$ as in
\eqref{eq:lorentzian_dirac} and $\Krein_{\nmax,\Nt} =
\HGV^{\nmax} \otimes \C^{\Nt}$,
\begin{equation}
\Lprop(\Tcal^{L}_{\nmax,\Nt,T}, \Tcal^{L}_{\Manifold})
\;\le\;
C_{3}^{\mathrm{joint}} \cdot
\gamma^{\mathrm{joint}}_{\nmax,\Nt,T}
\;\longrightarrow\; 0
\quad \text{as } (\nmax, \Nt) \to (\infty, \infty),
\label{eq:p45_main}
\end{equation}
with $C_{3}^{\mathrm{joint}} \le 1$ asymptotically tight (inherited
from Paper~38 L3 via the joint Lichnerowicz structural identity) and
\begin{equation}
\gamma^{\mathrm{joint}}_{\nmax,\Nt,T}
\;=\;
O\!\bigl(\log \nmax / \nmax \;+\; T / \Nt\bigr).
\label{eq:p45_rate}
\end{equation}
\end{theorem}

\subsection{Three structural simplifications}

The proof of Theorem~\ref{thm:p45_main} transfers the Paper~38
five-lemma chain to the joint $\SU(2) \times \Uone$ setting.  Three
structural simplifications arise that are not visible in the
Riemannian setting, and are recorded as substantive new content of
the Paper~45 proof:

\paragraph{PURE\_TENSOR Berezin map.}  The joint Berezin map
$B^{\mathrm{joint}} = B^{\SU(2)} \otimes B^{\Uone}$ factorises as a
pure tensor.  Each of the four required properties (positivity,
contractivity, approximate identity, L3 compatibility) inherits
factor-by-factor; no cross-factor verification is needed.

\paragraph{Vanishing time-chirality cross-term.}  The joint
Lichnerowicz commutator simplifies via the structural identity
\begin{equation}
[\DL,\, a_{s} \otimes a_{t}]
\;=\;
i\,[\DGV, a_{s}] \otimes a_{t}
\label{eq:vanishing_cross_term}
\end{equation}
because $\{\gamma^{0}, \DGV\} = 0$ on the chirality-doubled basis.
($\gamma^{0}$ anticommutes with the spatial Camporesi--Higuchi Dirac
because $\DGV$ is chirality-diagonal in the Peskin--Schroeder chiral
basis.)  The temporal direction contributes nothing to the joint
commutator, and $C_{3}^{\mathrm{joint}} = C_{3}^{\SU(2)}$ verbatim.

\paragraph{$\Nt$-independent joint cb-norm.}
By the Bo\.zejko--Fendler central-multiplier equality
\cite{bozejko_fendler1991} on the amenable group product
$\SU(2) \times \Uone$, the joint cb-norm of the central spectral
Fej\'er multiplier is
\begin{equation}
\cbnorm{S_{K^{\mathrm{joint}}}}
\;=\;
\frac{2}{\nmax + 1},
\label{eq:joint_cb_norm}
\end{equation}
\emph{independent} of $\Nt$.  This reflects the trivial
central-multiplier structure of the compact abelian factor $\Uone$,
and means that the $\Nt$-dependence in
\eqref{eq:p45_rate} lives entirely in the joint mass-concentration
moment (the $T/\Nt$ Fej\'er-on-the-circle rate), not in the
cb-norm height bound.

\paragraph{Riemannian-limit recovery at $\Nt = 1$ bit-exact.}
At $\Nt = 1$, the joint Berezin reduces bit-identically to
$B^{\SU(2)} \otimes B^{\Uone}_{1,T} = B^{\SU(2)}$, since the
$\Nt = 1$ temporal Berezin map is the identity on the single Fourier
mode.  Hence $\Lprop(\Tcal^{L}_{\nmax,1,T}, \Tcal^{L}_{\Manifold})$
agrees with the Riemannian Paper~38 bound bit-exactly, providing the
load-bearing falsifier check for the joint construction.

\subsection{Honest scope}

Theorem~\ref{thm:p45_main} is, to the author's knowledge, the first
published Lorentzian propinquity convergence theorem in the math.OA
literature.  Several aspects of its scope require explicit comment.

\begin{itemize}\setlength\itemsep{0pt}
\item[(Q1)] \textbf{Strong-form Lorentzian propinquity (Paper~45
\S 1.4 G1).}
The K$^{+}$-restriction reduces the indefinite Krein product to a
Hilbert inner product on $\Kplus$ and lets the standard Latr\'emoli\`ere
machinery apply verbatim.  Strong-form Lorentzian propinquity
(Latr\'emoli\`ere-style metric on Krein spectral triples in their own
right, without K$^{+}$-restriction) requires an indefinite-inner-product
analog of the operator-norm Lipschitz seminorm with appropriate
behaviour under Krein-self-adjoint $\DL$.  This remains an open
NCG-math problem; multi-month scale.
\item[(G2)] \textbf{De-compactification $T \to \infty$.}  The temporal
factor is compactified to $\Uone_{T}$ with finite radius $T$.
De-compactification to non-compact $\R_{t}$ is a separate program
(internal Sprint L3c) requiring renormalisation of the cb-norm bound
and reach moment.
\item[(G3)] \textbf{Cross-manifold $\Tcal_{\sthree} \otimes
\Tcal_{\mathrm{Hardy}(\sfive)}$.}  Blocked at the NCG-framework level
by the Riemannian-vs-Hardy-sector asymmetry recorded in Paper~24~\S V
of the GeoVac corpus (a four-layer Coulomb/HO asymmetry); this is
the Coulomb-side claim only.
\item[(G4)] \textbf{Inner-factor calibration data.}  Higgs / Yukawa
selection question of the Connes--Marcolli SM construction;
orthogonal to the present convergence theorem, addressed by
Paper~32~\S 8.C verdict POSITIVE-THIN (construction admits Higgs
structurally, but no autonomous Yukawa selection from the present
data).
\end{itemize}

% =====================================================================
\section{The master Mellin engine and the cross-arc connection}
\label{sec:mellin}
% =====================================================================

The most striking structural observation that emerges from reading
the arc is that the same transcendental signatures appear in
\emph{different} settings.  The constant $4/\pi$ of Paper~38 and the
constant $1/\pi^{2}$ of Paper~43~\S 10.2 are not numerical
coincidences: they are two appearances of the same Hopf-base measure
mechanism, the $M_{1}$ sub-case of the master Mellin engine
catalogued in Paper~32~\S 8.

\subsection{The $\pi$-source case-exhaustion theorem (Paper~32~\S 8)}

\begin{theorem}[Paper~32~\S 8, $\pi$-source case-exhaustion]
\label{thm:pi_source}
Let $\mathcal{C}$ be any finite composition of GeoVac projections
applied to a Layer-1 datum on $(\AGV, \HGV, \DGV)$ or its
Bargmann--Segal sibling triple.  The following are equivalent:
\begin{itemize}\setlength\itemsep{0pt}
\item[(i)] $\mathcal{C}$(datum) contains a transcendental of the form
$\pi^{k}$ for $k \in \Q$;
\item[(ii)] $\mathcal{C}$ contains at least one projection whose
evaluation includes a continuous integration over a parameter
promoted from the discrete graph spectrum;
\item[(iii)] $\mathcal{C}$ engages at least one of three
spectral-triple mechanisms,
\begin{itemize}\setlength\itemsep{0pt}
\item[$M_{1}$] (\emph{Hopf-base measure}):\
$\pi = \Vol(S^{2})/4$ via the Hopf bundle $\sthree \to S^{2} \times
S^{1}$, equivalently $4/\pi = \Vol(S^{2})/\pi^{2}$;
\item[$M_{2}$] (\emph{Seeley--DeWitt}):\
$\sqrt{\pi}$ via heat-kernel coefficients on the Camporesi--Higuchi
Dirac spectrum on $\sthree$ (Paper~28 Theorem~T9:
$\zeta_{D^{2}}(s)$ is polynomial in $\pi^{2}$ with rational
coefficients at every integer $s$);
\item[$M_{3}$] (\emph{vertex-topology Dirichlet $L$}):\
Catalan $G$ and $\beta(s)$ via quarter-integer Hurwitz shifts in
two-loop sunset vertices.
\end{itemize}
\end{itemize}
\end{theorem}

The TS-E3 sharpening of Paper~32 further shows that $M_{1}, M_{2},
M_{3}$ are sub-cases of a single master Mellin-engine mechanism
\begin{equation}
\pi\text{-source}
\;=\;
\mathcal{M}\bigl[ \Tr(D^{k} e^{-tD^{2}}) \bigr]
\quad \text{with } k \in \{0, 1, 2\},
\label{eq:master_mellin_engine}
\end{equation}
where $k = 0$ selects $M_{1}$, $k = 1$ selects $M_{3}$, and $k = 2$
selects $M_{2}$.

\subsection{$M_{1}$ in the Riemannian GH rate}

The rate constant $4/\pi$ in \eqref{eq:gh_rate_constant} is the
$M_{1}$ Hopf-base measure signature.  This is not a derived
observation; it is forced by the Plancherel-weight $\times$
Vandermonde-Jacobian cancellation theorem (Paper~40~\S 3.2) which
leaves only the rank-1 Stein--Weiss constant after the
$\prod_{\alpha > 0}$ products cancel.  Equivalently, the
single-positive-root Euler--Catalan sum
\begin{equation*}
\sum_{d \text{ odd}} \frac{1}{d^{2}} \;=\; \frac{\pi^{2}}{8}
\end{equation*}
produces the $\pi^{2}$ in the denominator, and the Hopf bundle
$\sthree \to S^{2}$ produces $\Vol(S^{2}) = 4\pi$ in the numerator.

\subsection{$M_{1}$ in the Lorentzian HS-orthogonality}

The $1/\pi^{2}$ prefactor on $\|H_{\mathrm{local}}\|^{2}$ in the
Pythagorean closed form \eqref{eq:hs_closed_form} is the same
$M_{1}$ signature \emph{at the operator-system level}: the
Krein wedge BW choice $H_{\mathrm{local}} = K_{\alpha}^{W}/(2\pi)$
introduces a factor of $1/(2\pi)$ from the modular-circle
circumference, which squared gives $1/(4\pi^{2})$ in
$\|H_{\mathrm{local}}\|^{2}$.  The Hopf-base measure factor $4/\pi$
of Paper~38 and the $1/\pi^{2}$ of Paper~43 are two manifestations
of the same $\Vol(S^{2}) = 4\pi$ identity in the master Mellin
engine.

\paragraph{Why this matters.}
The case-exhaustion theorem (Theorem~\ref{thm:pi_source}) is a
\emph{spectral-triple-level} statement about which transcendentals
can appear in evaluations of GeoVac projection chains.  The constant
$4/\pi$ of Paper~38 is a \emph{metric-geometry-level} statement about
the propinquity rate.  The fact that these two abstractly-different
levels agree on the same signature is non-trivial: it identifies the
metric geometry of the GeoVac spectral triple as carrying the same
M$_{1}$ signature that the abstract Mellin engine catalogues.
Abstract taxonomy meets concrete propinquity rates.  Paper~43~\S 10.2
takes this further: the Lorentzian Pythagorean decomposition at the
operator-system level carries the same M$_{1}$ signature again, this
time on the operator algebra of the Krein wedge.

% =====================================================================
\section{Open questions and outlook}
\label{sec:open}
% =====================================================================

The arc closes the convergence side of the operator-algebra story at
the K$^{+}$-weak-form level (Paper~45) and the modular-Hamiltonian
side at the operator-system level on Riemannian (Paper~42) and
Lorentzian (Paper~43) backgrounds.  Several questions remain open.

\paragraph{Strong-form Lorentzian propinquity (Q1).}
As noted in \S\ref{sec:lorentzian}, the K$^{+}$-restriction is the
move that lets the standard Latr\'emoli\`ere machinery apply
verbatim.  Strong-form Lorentzian propinquity, formulated as a
Latr\'emoli\`ere-style metric on Krein spectral triples in their
own right, would require an indefinite-inner-product analog of the
operator-norm Lipschitz seminorm with the right behaviour under
Krein-self-adjoint $\DL$.  This is a genuine NCG-math problem at
multi-month scale.  Paper~45~\S 1.4 G1 names this explicitly.

\paragraph{De-compactification $T \to \infty$ (G2).}
The temporal factor $\Uone_{T}$ is finite; the standard physics
applications (Hawking, Unruh) require non-compact $\R_{t}$.  The
rate $T/\Nt \to 0$ in \eqref{eq:p45_rate} is the joint
mass-concentration on $\Uone_{T}$ as $\Nt \to \infty$ at fixed $T$;
the $T \to \infty$ limit is a separate program (Sprint L3c).

\paragraph{Cross-manifold $\Tcal_{\sthree} \otimes
\Tcal_{\mathrm{Hardy}(\sfive)}$ (G3, Coulomb-side claim only).}
The natural tensor product of the Coulomb $\sthree$ spectral triple
with the Bargmann--Segal Hardy-sector spectral triple on $\sfive$
\cite{paper24} is blocked at the NCG-framework level by a four-layer
asymmetry recorded in Paper~24~\S V.  The four layers are:
(i) spectrum-computing role of $L_{0}$ (Coulomb-side only, by Paper~23
Fock rigidity);
(ii) calibration $\pi$ (Coulomb-side only);
(iii) non-abelian Wilson gauge with natural matter coupling (both
sides have Wilson, only Coulomb-side has natural matter coupling
because the Bargmann tower $(N,0)$ irreps have $N$-dependent
dimensions and no uniform-rep matter sector); and (iv) modular-
Hamiltonian structure of the wedge KMS state (Coulomb-side admits the
Krein wedge HS-orthogonality of Corollary~\ref{cor:hs_orthogonality}
with closed-form $M_{1}$ prefactor; the Hardy-sector side does not
admit the construction at all).  Closing G3 requires new
NCG-framework machinery to bridge Riemannian and Hardy-sector spectral
triples in a single tensor product; we make no claim about whether
this is achievable.

\paragraph{Higgs as inner fluctuation (G4).}
The GeoVac triple is naturally combined with a finite almost-
commutative factor $\AGV \otimes (\C \oplus \mathbb{H})$ to produce
an electroweak-like sector.  The Connes axiom audit at finite cutoff
(Paper~32~\S 8.C) gives verdict POSITIVE-THIN: the construction
\emph{admits} a Higgs sector structurally (i.e.~the inner fluctuation
$D \mapsto D + \omega + \varepsilon' J \omega J^{-1}$ produces a
non-trivial off-diagonal $\mathbb{C} \leftrightarrow \mathbb{H}$
block iff the finite Yukawa $Y_{F}$ is non-zero), but \emph{GeoVac
does not autonomously select} $Y_{F}$ from its internal data.  In
Marcolli--vS terms, GeoVac sits on the Yang--Mills-without-Higgs side
of the Perez-Sanchez \cite{perez_sanchez2024,perez_sanchez2025}
clarification, with calibration data (the Yukawa) supplied externally.
Paper~32~\S 8.C records this as the standing structural verdict.

\paragraph{State-side complement of the dictionary.}
Paper~44 highlights that the non-trivial K$^{+}$-positivity content of
the Lorentzian operator system sits at the \emph{state} level, not the
operator level.  The natural object is the Wasserstein--Kantorovich
metric on $K^{+}$-states.  Paper~45's K$^{+}$-weak-form propinquity is
the first concrete construction in this direction; substantially more
mathematics on Wasserstein--Kantorovich on Krein-positive states is
required to complete the state-side dictionary.

\paragraph{Concluding remark.}
The structural identification of GeoVac as an almost-commutative
spectral triple in the Marcolli--van~Suijlekom lineage is proved at
the GH-convergence level by Paper~38 (Riemannian, single factor),
extended by Paper~39 (tensor product), generalised by Paper~40
(all compact simple Lie groups), and extended in turn to the
Lorentzian K$^{+}$-weak-form by Paper~45.  The metric geometry of
the truncated triples carries the same M$_{1}$ Hopf-base measure
signature ($4/\pi$ in the Riemannian rate, $1/\pi^{2}$ in the
Lorentzian Pythagorean closed form) that the abstract
$\pi$-source case-exhaustion theorem (Theorem~\ref{thm:pi_source})
identifies at the spectral-triple level.  The cross-arc
identification of these signatures is the deepest structural
observation of the corpus.

% =====================================================================
% Bibliography
% =====================================================================

\bibliographystyle{plainnat}
\begin{thebibliography}{99}

\bibitem[Bisognano \& Wichmann(1976)]{bisognano_wichmann1976}
J.~J.~Bisognano and E.~H.~Wichmann,
\newblock \emph{On the duality condition for quantum fields},
\newblock J.~Math.~Phys. \textbf{17} (1976), 303--321.

\bibitem[Bizi, Brouder \& Besnard(2018)]{bizi_brouder_besnard2018}
N.~Bizi, C.~Brouder, and F.~Besnard,
\newblock \emph{Space and time dimensions of algebras with
application to Lorentzian noncommutative geometry and quantum
electrodynamics},
\newblock J.~Math.~Phys. \textbf{59} (2018), 062303;
\newblock arXiv:1611.07062.

\bibitem[Bo\.zejko \& Fendler(1991)]{bozejko_fendler1991}
M.~Bo\.zejko and G.~Fendler,
\newblock \emph{Herz--Schur multipliers and uniformly bounded
representations of discrete groups},
\newblock Arch.~Math. \textbf{57} (1991), 290--298.

\bibitem[Camporesi \& Higuchi(1996)]{camporesi_higuchi1996}
R.~Camporesi and A.~Higuchi,
\newblock \emph{On the eigenfunctions of the Dirac operator on
spheres and real hyperbolic spaces},
\newblock J.~Geom.~Phys. \textbf{20} (1996), 1--18;
\newblock arXiv:gr-qc/9505009.

\bibitem[Chamseddine \& Connes(2010)]{chamseddine_connes2010}
A.~H.~Chamseddine and A.~Connes,
\newblock \emph{Noncommutative geometry as a framework for
unification of all fundamental interactions including gravity. I},
\newblock Fortsch.~Phys. \textbf{58} (2010), 553--600.

\bibitem[Connes(1995)]{connes1995}
A.~Connes,
\newblock \emph{Noncommutative Geometry},
\newblock Academic Press, 1995.

\bibitem[Connes \& van~Suijlekom(2021)]{connes_vs2021}
A.~Connes and W.~D.~van~Suijlekom,
\newblock \emph{Spectral truncations in noncommutative geometry and
operator systems},
\newblock Comm.~Math.~Phys. \textbf{383} (2021), 2021--2067;
\newblock arXiv:2004.14115.

\bibitem[Farsi \& Latr\'emoli\`ere(2024)]{farsi_latremoliere2024}
C.~Farsi and F.~Latr\'emoli\`ere,
\newblock \emph{Continuity of spectral propinquity for spectrally
truncated spheres},
\newblock preprint, 2024.

\bibitem[Farsi \& Latr\'emoli\`ere(2025)]{farsi_latremoliere2025}
C.~Farsi and F.~Latr\'emoli\`ere,
\newblock \emph{Continuity for the spectral propinquity of the
Dirac operators associated with an analytic path of Riemannian
metrics},
\newblock arXiv:2504.11715, 2025.

\bibitem[Hartle \& Hawking(1976)]{hartle_hawking1976}
J.~B.~Hartle and S.~W.~Hawking,
\newblock \emph{Path-integral derivation of black-hole radiance},
\newblock Phys.~Rev. D \textbf{13} (1976), 2188--2203.

\bibitem[Hawkins(2000)]{hawkins2000}
E.~Hawkins,
\newblock \emph{Quantization of equivariant vector bundles, and the
correspondence with deformation quantization},
\newblock Comm.~Math.~Phys. \textbf{215} (2000), 409--432.

\bibitem[Hekkelman(2022)]{hekkelman2022}
E.~Hekkelman,
\newblock \emph{Truncations of the circle and Connes' geometric
formula},
\newblock M.Sc.\ thesis, Radboud University, 2022;
\newblock arXiv:2206.13744.

\bibitem[Hekkelman \& McDonald(2024)]{hekkelman_mcdonald2024}
E.~Hekkelman and E.~McDonald,
\newblock \emph{Spectral truncations of $\T^{d}$ and quantum metric
geometry},
\newblock J.~Noncommut.~Geom., to appear (2024);
\newblock arXiv:2403.18619.

\bibitem[Hekkelman \& McDonald(2024b)]{hekkelman_mcdonald2024b}
E.~Hekkelman and E.~McDonald,
\newblock \emph{A noncommutative integral on spectrally truncated
spectral triples, and a link with quantum ergodicity},
\newblock J.~Funct.~Anal., to appear (2025);
\newblock arXiv:2412.00628.

\bibitem[Latr\'emoli\`ere(2018/2023)]{latremoliere_metric_st_2017}
F.~Latr\'emoli\`ere,
\newblock \emph{The Gromov--Hausdorff propinquity for metric
spectral triples},
\newblock Adv.~Math. \textbf{415} (2023), Paper No.~108876, 88pp;
\newblock preprint arXiv:1811.10843, November 2018.

\bibitem[Latr\'emoli\`ere(2018)]{latremoliere2018}
F.~Latr\'emoli\`ere,
\newblock \emph{The Gromov--Hausdorff propinquity},
\newblock Trans.~Amer.~Math.~Soc. \textbf{370} (2018), 365--411.

\bibitem[Leimbach \& van~Suijlekom(2024)]{leimbach_vs2024}
L.~Leimbach and W.~D.~van~Suijlekom,
\newblock \emph{Gromov--Hausdorff convergence of spectral
truncations of the torus},
\newblock Adv.~Math. \textbf{439} (2024), Paper No.~109496.

\bibitem[Marcolli \& van~Suijlekom(2014)]{marcolli_vs2014}
M.~Marcolli and W.~D.~van~Suijlekom,
\newblock \emph{Gauge networks in noncommutative geometry},
\newblock J.~Geom.~Phys. \textbf{75} (2014), 71--91;
\newblock arXiv:1301.3480.

\bibitem[Mondino \& S\"amann(2024)]{mondino_samann2024}
A.~Mondino and C.~S\"amann,
\newblock \emph{Lorentzian metric measure spaces},
\newblock Lett.~Math.~Phys. \textbf{114} (2024), Paper No.~37.

\bibitem[Mondino \& S\"amann(2025)]{mondino_samann2025}
A.~Mondino and C.~S\"amann,
\newblock \emph{Synthetic Lorentzian Gromov--Hausdorff convergence
and pre-compactness},
\newblock arXiv:2504.10380, 2025.

\bibitem[Paulsen(2002)]{paulsen2002}
V.~I.~Paulsen,
\newblock \emph{Completely Bounded Maps and Operator Algebras},
\newblock Cambridge Studies in Advanced Mathematics 78,
Cambridge University Press, 2002.

\bibitem[Perez-Sanchez(2024)]{perez_sanchez2024}
C.~I.~Perez-Sanchez,
\newblock \emph{On the continuum limit of gauge networks},
\newblock arXiv:2401.03705, 2024.

\bibitem[Perez-Sanchez(2025)]{perez_sanchez2025}
C.~I.~Perez-Sanchez,
\newblock \emph{Yang--Mills theories from gauge networks without
Higgs},
\newblock arXiv:2508.17338, 2025.

\bibitem[Pier(1984)]{pier1984}
J.-P.~Pier,
\newblock \emph{Amenable Locally Compact Groups},
\newblock Wiley, 1984.

\bibitem[Pisier(2001)]{pisier2001}
G.~Pisier,
\newblock \emph{Similarity Problems and Completely Bounded Maps},
\newblock Lecture Notes in Math.~1618, Springer, 2nd ed., 2001.

\bibitem[Sewell(1982)]{sewell1982}
G.~L.~Sewell,
\newblock \emph{Quantum fields on manifolds:\ PCT and gravitationally
induced thermal states},
\newblock Ann.~Phys.\ (NY) \textbf{141} (1982), 201--224.

\bibitem[Stein \& Weiss(1971)]{stein_weiss1971}
E.~M.~Stein and G.~Weiss,
\newblock \emph{Introduction to Fourier Analysis on Euclidean
Spaces},
\newblock Princeton Math.~Ser.~32, Princeton Univ.~Press, 1971.

\bibitem[Toyota(2023)]{toyota2023}
M.~Toyota,
\newblock \emph{Quantum Gromov--Hausdorff convergence of spectral
truncations for groups with polynomial growth},
\newblock arXiv:2309.13469, 2023.

\bibitem[Unruh(1976)]{unruh1976}
W.~G.~Unruh,
\newblock \emph{Notes on black-hole evaporation},
\newblock Phys.~Rev.\ D \textbf{14} (1976), 870--892.

\bibitem[van~den~Dungen(2016)]{vandungen2016}
K.~van~den~Dungen,
\newblock \emph{Krein spectral triples and the fermionic action},
\newblock Math.~Phys.~Anal.~Geom. \textbf{19} (2016), Art.~4, 22pp;
\newblock arXiv:1505.01939.

% --- Internal GeoVac papers (Group 1 operator-algebra arc) ---

\bibitem[Loutey, internal Paper~24]{paper24}
J.~Loutey,
\newblock \emph{Bargmann--Segal lattice for the 3D harmonic
oscillator and the Coulomb/HO structural asymmetry},
\newblock GeoVac internal preprint (2026), DOI via Zenodo.

\bibitem[Loutey, internal Paper~29]{paper29}
J.~Loutey,
\newblock \emph{The GeoVac Hopf graphs are Ramanujan:\ Ihara zeta,
integer-algebraic zeros, and the scope boundary on Selberg-on-
hydrogen},
\newblock GeoVac internal preprint (2026), DOI via Zenodo.

\bibitem[Loutey, internal Paper~32]{paper32}
J.~Loutey,
\newblock \emph{The GeoVac spectral triple, the master Mellin engine
case-exhaustion theorem, and the GH-convergence theorem on
$\sthree$},
\newblock GeoVac internal preprint (2026), DOI via Zenodo.

\bibitem[Loutey, internal Paper~38]{paper38}
J.~Loutey,
\newblock \emph{Latr\'emoli\`ere propinquity convergence of truncated
Camporesi--Higuchi spectral triples on $\sthree$},
\newblock GeoVac internal preprint (2026), DOI via Zenodo.

\bibitem[Loutey, internal Paper~39]{paper39}
J.~Loutey,
\newblock \emph{Tensor-product propinquity convergence for
distinct-focal-length truncated Camporesi--Higuchi spectral triples
on $\sthree$},
\newblock GeoVac internal preprint (2026), DOI via Zenodo.

\bibitem[Loutey, internal Paper~40]{paper40_unified}
J.~Loutey,
\newblock \emph{Latr\'emoli\`ere propinquity convergence of spectral
truncations of compact Lie groups with bi-invariant metric},
\newblock GeoVac internal preprint (2026), DOI via Zenodo.

\bibitem[Loutey, internal Paper~42]{paper42}
J.~Loutey,
\newblock \emph{Tomita--Takesaki modular structure on truncated
$\SU(2)$ spectral triples:\ four-witness Wick-rotation literal
identification at finite cutoff},
\newblock GeoVac internal preprint (2026), DOI via Zenodo.

\bibitem[Loutey, internal Paper~43]{paper43}
J.~Loutey,
\newblock \emph{Lorentzian extension of the four-witness
Wick-rotation theorem at finite cutoff},
\newblock GeoVac internal preprint (2026), DOI via Zenodo.

\bibitem[Loutey, internal Paper~44]{paper44}
J.~Loutey,
\newblock \emph{Operator-system extension of the BBB Krein spectral
triple at finite cutoff:\ propagation number and envelope dependence
on the Camporesi--Higuchi spinor bundle},
\newblock GeoVac internal preprint (2026), DOI via Zenodo.

\bibitem[Loutey, internal Paper~45]{paper45}
J.~Loutey,
\newblock \emph{K$^{+}$-restricted weak-form Lorentzian propinquity
convergence on truncated $\SU(2) \times \Uone_{T}$ Krein spectral
triples},
\newblock GeoVac internal preprint (2026), DOI via Zenodo.

\end{thebibliography}

\end{document}
